\emph{The finite set $\mathcal P$.} Define $\mathcal P$ as the set of the first $1000$ primes $\ell>10^9$ with $\ell\equiv65\pmod{128}$, i.e.\ the $1000$ smallest such primes. This is a finite, explicitly enumerable set: the primes $\equiv65\pmod{128}$ above $10^9$ form an increasing sequence, and $\mathcal P$ is its initial segment of length $1000$. Its largest element is $1001287361$, so $\mathcal P=\{\ell\text{ prime}:10^9<\ell\le1001287361,\ \ell\equiv65\pmod{128}\}$; the two descriptions agree because the second set has exactly $1000$ elements and contains the first. The endpoints coincide with those printed in \cite{KY} \S4.2.

\emph{Target-list certificate (claim KY1000\_TARGET, label C).} The shipped list (\path{sage/family_ky1000/KY1000_primes.txt}) is checked, by a script with no dependence on the generator, to have $1000$ entries, distinct and ascending, every entry prime (deterministic primality test for this size), every entry $\equiv65\pmod{128}$ and $>10^9$; the script then regenerates the first $1000$ qualifying primes from $10^9$ by a sieve and compares elementwise, and finally compares the list with the set of primes in the witness-file names. The target-list checker also rejects five planted corruptions of the list (one entry altered to a composite in the wrong class; one qualifying prime omitted and the list padded with the $1001$st; the last entry duplicated; two neighbours swapped; one prime replaced by the next qualifying prime beyond $1001287361$, which only the interval form catches; verifier steps 03b/03c, a control set separate from the nine witness-file controls). This discharges the identification of the shipped list with the set $\mathcal P$ defined above.

\emph{Witness certificates (Lemma~\ref{lem:ky1000}, label C).} For every $\ell\in\mathcal P$ there is one witness file (SHA-256 in the manifest). The read-only verifier \path{scripts/family_verify.sage} (specification in the paper's certificate section; frozen; SHA-256 \texttt{0f38bb0d}\ldots) accepts each of them with verdict EXCLUDED: the verification ledger records EXCLUDED $1000/1000$, every one of the $32=2^{7-2}$ components of every prime carrying an accepted T certificate ($32000$ T lines) and $31987$ of them also an accepted RHO line; the largest accepted $T$ value is $4164.4897<\bar T_7=4224$. Nine planted corruptions of a passing file (wrong factor, wrong vector, vector in $\ell R_n$, altered threshold, altered header fields, \ldots) are REJECTED and the untouched $n=2$ file is accepted (\path{sage/negctl_r13/negctl_r13.log}; ledger \path{certificates/negctl/negctl_ledger_r14.json}).

\emph{From EXCLUDED to $\ell\nmid h_7/h_6$ (Theorem~\ref{thm:cert}, C+L, F-core).} By Theorem~\ref{thm:cert}, EXCLUDED for $(7,\ell)$ implies $\ell\nmid k_7=h_7/h_6$: each accepted line is a nonzero vector of the component lattice refuting saturation of that component (RHO: height below the floor $L_7$; T: trace of the square below $\bar T_7$, \cite{KY} Thm.~2.3), and saturation of some component is forced by $\ell\mid k_7$ (\cite{KY} Prop.~4.1). Both refutations are instances of the same discrete core as Theorem~A (\path{WeberCert.exclusion_chain_direct}); the literature inputs are those of the trust table of the paper (KY Eq.~(17), KY Prop.~4.1, KY Def.~2.1, MO16 Lemma~2.5(1), Blichfeldt) plus the trace floor KY Thm.~2.3 / MO3 Prop.~6.6.

\emph{Scope.} Nothing is claimed for primes outside $\mathcal P$, for other residue classes, or for $n\ne7$; \cite{KY} had excluded these $1000$ primes conditionally on their Conjecture~2.2, which is not used here.
